\documentclass[10pt,a4paper]{amsart}
\usepackage[margin=19mm]{geometry}
\usepackage{amsmath,amssymb,mathtools}
\usepackage[scr=rsfs,cal=euler]{mathalfa}
\usepackage{xfrac}
\usepackage[unicode,hidelinks,bookmarks=false]{hyperref}
\newcommand{\ie}{i.e.\ }
\newcommand{\cf}{cf.\ }
\newcommand{\resp}{respectively\ }
\newcommand{\Subsection}{\subsection}
\providecommand{\nobreakdash}{\nobreak\hbox{-}}
\setlength{\emergencystretch}{2em}
\multlinegap0pt
\usepackage{amssymb}
\usepackage{bm}
\usepackage{xcolor}
\usepackage{enumerate}
\usepackage{tikz}
\newtheorem{proposition}{Proposition}
\renewcommand{\Re}{\operatorname{Re}}
\newcommand{\dd}{\mathrm{d}}
\title{Extreme central values of quadratic Dirichlet $L$-functions with prime conductors}
\author{Mingyue Fan, Shenghao Hua, Sizhe Xie}
\date{Source: arXiv:2306.16886v2 (CC BY 4.0)}
\begin{document}
\maketitle

\noindent\emph{Source and licence.} Extreme central values of quadratic Dirichlet $L$-functions with prime conductors. Version arXiv:2306.16886v2 (CC BY 4.0), distributed by arXiv under the Creative Commons Attribution 4.0 International licence, http://creativecommons.org/licenses/by/4.0/, as recorded in the arXiv licence metadata for this exact version and shown on the abstract page https://arxiv.org/abs/2306.16886v2. Full licence notice: \texttt{licenses/arxiv-2306.16886-CC-BY-4.0.txt}.\par\smallskip

\noindent\textit{Editorial scope.} Counted unit: the article's proposition prop:Mi (source0.tex lines 399--411). The article's complete original proof of it is delivered with it (source0.tex lines 436--657, one \texttt{proof} environment of 222 lines, which the article places away from the statement). No mathematics is written, repaired or re-derived: the statement and the proof are the article's own lines, in the article's own notation, and no retained line is substituted or re-flowed.  No further source-local context is needed: every cross-reference of every retained line resolves inside the two delivered spans.  Nothing was omitted from any retained span, so the package carries no omission record.  No retained line contains a citation, so the wrapper carries no bibliography and no bibliography entry.  No retained line contains a figure environment, an image-inclusion command or drawing code, so no image is delivered and the package declares no asset dependency.  Editorial formatting only, in addition: an AMS \texttt{amsart} preamble that restates the article's own declarations and macros used by the retained lines, namely the environments \texttt{proposition} on separate counters, together with the standard packages those lines need.  The retained environments print in this document as Proposition 1 in order of appearance, whereas the article prints the same items with the numbering of its own full document.  The complete native source of the article as distributed in the arXiv e-print for this exact version is bundled unchanged as \texttt{sources/143580/source0.tex}.
\begin{proposition}\label{prop:Mi}
For $\theta< 1/6$,
  \begin{equation*}
  \mathcal{M}_1\ll X\prod_{\mathcal{L}^2\leq p \leq \exp((\log \mathcal{L})^2)}
(1+ b(p)^2)
\end{equation*}
and
\begin{equation*}
   |\mathcal{M}_2|\gg
   X\log X \prod_{\mathcal{L}^2\leq p \leq \exp((\log \mathcal{L})^2)}
  \bigg(1+2\frac{b(p)}{\sqrt{p}}+b(p)^2\bigg).
\end{equation*}
\end{proposition}

\begin{proof}[Proof of Proposition \ref{prop:Mi}]
For $m_1,m_2$ square-free, we write $m_1m_2=m_0(m_1,m_2)^2$ so that $m_0$ is square-free.
Thus
\begin{equation*}
\begin{split}
  \mathcal{M}_2&=
   \frac{2}{(1-\frac{1}{\sqrt{2}})^{2}}
  \underset{m_1,m_2\leq M}{\sum\sum}
  b(m_1)b(m_2)
  \sum_{p\equiv 1\pmod 8 }
  (\log p)\Phi\bigg(\frac{p}{X}\bigg)
  \\
&\hskip 30pt\times
  \frac{1}{2\pi i}
  \int_{(u)}
\frac{\Gamma(\frac{s}{2}+\frac{1}{4})}
{\Gamma(\frac{1}{4})}
(1-2^{-\frac{1}{2}+s})
  \sum_{\substack{nm_0=\square\\(n,2p)=1}}
\frac{1}
{n^{\frac{1}{2}+s}}
\bigg(\frac{p}{\pi}\bigg)^{\frac{s}{2}}
e^{s^2}\frac{\dd s}{s}\\
&=
 \frac{2}{(1-\frac{1}{\sqrt{2}})^{2}}
 \underset{m_1,m_2\leq M}{\sum\sum}
  \frac{b(m_1)b(m_2)}{\sqrt{m_0}}
  \sum_{p\equiv 1\pmod 8 }
  (\log p)\Phi\bigg(\frac{p}{X}\bigg)
  \\
&\hskip 30pt\times
  \frac{1}{2\pi i}
  \int_{(u)}
\frac{\Gamma(\frac{s}{2}+\frac{1}{4})}
{\Gamma(\frac{1}{4})}
\zeta(1+2s)\bigg(\frac{p}{\pi m_0^2}\bigg)^{\frac{s}{2}}
(1-2^{-\frac{1}{2}+s})
(1-2^{-1-2s})
(1-p^{-1-2s})
e^{s^2}\frac{\dd s}{s},
\end{split}
\end{equation*}

Let $s=u+it$, and
\begin{equation*}
  \mathcal{I}_{p,m_0}=\frac{1}{2\pi i}
  \int_{(u)}
\frac{\Gamma(\frac{s}{2}+\frac{1}{4})}
{\Gamma(\frac{1}{4})}
\zeta(1+2s)
\bigg(\frac{p}{\pi m_0^2}\bigg)^{\frac{s}{2}}
(1-2^{-\frac{1}{2}+s})
(1-2^{-1-2s})
(1-p^{-1-2s})
e^{s^2}\frac{\dd s}{s},
\end{equation*}
we move the line of integration to $\Re(s)=-\frac{1}{2}+\varepsilon$,
leaving a residue at $s=0$,
the new integral is
\begin{equation*}
\begin{split}
  &\frac{1}{2 \pi^{3/4 + \varepsilon/2}}
  p^{-\frac{1}{4}+\frac{\varepsilon}{2}} m_0^{\frac{1}{2}-\varepsilon}\\
&\hskip 20pt\times
\int_{-\infty}^{+\infty}
\frac{\Gamma(\frac{\varepsilon+it}{2})}
{\Gamma(\frac{1}{4})}
\zeta(2\varepsilon+2it)
\bigg(\frac{p}{\pi m_0^2}\bigg)^{\frac{it}{2}}
(1-2^{-1+\varepsilon+it})
(1-2^{-2\varepsilon-2it})
(1-p^{-2\varepsilon-2it})
\frac{e^{(-\frac{1}{2}+\varepsilon+it)^2}}
{-\frac{1}{2}+\varepsilon+it}\dd t\\
&\ll p^{-\frac{1}{4}+\varepsilon}
m_0^{\frac{1}{2}}
\int_{-\infty}^{+\infty}
|\zeta(2\varepsilon+2it)|
\frac{e^{-t^2}}{1+|t|} \dd t\\
&\ll p^{-\frac{1}{4}+\varepsilon}
m_0^{\frac{1}{2}}.
\end{split}
\end{equation*}
We have the Laurent series expansions in a neighborhood of $s=0$
\begin{equation*}
\begin{gathered}
\frac{\Gamma(\frac{s}{2}+\frac{1}{4})}
{\Gamma(\frac{1}{4})}
=1+\frac{1}{2}\frac{\Gamma^{'}(\frac{1}{4})}
{\Gamma(\frac{1}{4})}s+\dots
=1-\bigg(\frac{\pi}{4}+\frac{3\log 2}{2} +\frac{\gamma}{2}\bigg)s+\dots;\\
\zeta(1+2s)=\frac{1}{2s}+\gamma+\dots;\\
\bigg(\frac{p}{\pi m_0^2}\bigg)
^{\frac{s}{2}}=1+\frac{1}{2}
\log\bigg(\frac{p}{\pi m_0^2}\bigg)s+\dots;\\
(1-2^{-\frac{1}{2}+s})
(1-2^{-1-2s})
(1-p^{-1-2s})
e^{s^2}
=2^{-1}(1-2^{-\frac{1}{2}})
(1-p^{-1})
\\ \times
\bigg(1+\bigg(2+2\log 2+\frac{2\log p}{p-1}-\frac{\log 2}{\sqrt{2}-1}\bigg)s+\dots\bigg);
\end{gathered}
\end{equation*}
where $\gamma$ is the Euler constant.
It follows that the residues can be written as
\begin{equation}\label{eq:residue}
  2^{-1}(1-2^{-\frac{1}{2}})
(1-p^{-1})\Big(\frac{1}{4}\log\bigg(\frac{p}{m_0^2}\bigg)
+\frac{\log p}{p-1}+ c_0
\Big)
\end{equation}
where fixed constant
\begin{equation*}
  c_0=-\bigg(\frac{\pi}{4}+\frac{3\log 2}{2} +\frac{\gamma}{2}\bigg)
-\frac{\log(\pi)}{4}
+\frac{3}{2}-\frac{\log 2}{2\sqrt{2}-2}+\gamma,
\end{equation*}
and $c_0\approx -1.15939623$ is easily known by hand or by MATLAB.
Then for $X$ large enough, $p\geq X$, we have
\begin{equation*}
  |\mathcal{I}_{p,m_0}|\geq
  2^{-1}(1-2^{-\frac{1}{2}})
(1-X^{-1})\bigg(\frac{1}{4}
\log\bigg(\frac{p}{m_0^2}\bigg)
-1.15939624\bigg)
\gg \log X
\end{equation*}
since $m_0\le M^2=X^{2\theta}$.

Therefore, we can get
\begin{equation}\label{M2 gg}
\begin{split}
  |\mathcal{M}_2|&\gg
 \underset{m_1,m_2\leq M}{\sum\sum}
  \frac{b(m_1)b(m_2)}{\sqrt{m_0}}
  \sum_{p\equiv 1\pmod 8 }
  (\log p)\Phi\bigg(\frac{p}{X}\bigg)\log X\\
  &\gg_\Phi X\log X
  \underset{m_1,m_2\leq M}{\sum\sum}
  \frac{b(m_1)b(m_2)}{\sqrt{m_0}},
  \end{split}
\end{equation}
by applying the prime number theorem in arithmetic progressions
\begin{equation*}
  \sum_{\substack{p\equiv 1\pmod 8 \\ p \textrm{ prime }}}
  (\log p)\Phi\bigg(\frac{p}{X}\bigg)
  =\frac{X}{4}\check{\Phi}(0)+O\bigg(\frac{X}{\log X}\bigg),
\end{equation*}
where $m_0=\frac{m_1m_2}{(m_1, m_2)^2}$.

Both $m_1,m_2$ in \eqref{M2 gg} are square-free, then we write $m_1=qr$, $m_2=qs$, $m_0=rs$ with $(q,r)=(q,s)=(r,s)=1$, we have
\begin{equation}\label{truncated sum}
    \underset{m_1,m_2\leq M}{\sum\sum}
  \frac{b(m_1)b(m_2)}{\sqrt{m_0}}
  =\underset{\substack{(q,r)=(q,s)=(r,s)=1
  \\qr,qs\leq M}}
  {\sum\sum\sum}
  \frac{b(qr)b(qs)}{\sqrt{rs}}.
\end{equation}
We remove the restriction of $qr,qs\leq M$, the extend sum is
\begin{equation}\label{complete sum}
\underset{\substack{(q,r)=(q,s)=(r,s)=1}}
  {\sum\sum\sum}
  \frac{b(qr)b(qs)}{\sqrt{rs}}
    =\prod_{\mathcal{L}^2\leq p \leq \exp((\log \mathcal{L})^2)}
  \bigg(1+2\frac{b(p)}{\sqrt{p}}+b(p)^2\bigg).
\end{equation}

By the symmetry of $qr$ and $qs$, the difference between \eqref{truncated sum} and \eqref{complete sum} could be controlled by
  \begin{equation*}
2\underset{\substack{(q,r)=(q,s)=(r,s)=1
  \\qr> M}}
  {\sum\sum\sum}
  \frac{b(q)^2b(r)b(s)}{\sqrt{rs}}.
\end{equation*}
By employing Rankin's trick, for any $\alpha>0$, the above is
\begin{equation}\label{eqn:Rankinerror}
\begin{split}
 &\ll M^{-\alpha}\underset{
 \substack{(q,r)=(q,s)=(r,s)=1}}
  {\sum\sum\sum}b(q)^2 b(r)b(s)
  q^{\alpha}r^{-\frac{1}{2}+\alpha}
s^{-\frac{1}{2}}\\
&\ll M^{-\alpha}
\prod_{\mathcal{L}^2\leq p \leq \exp((\log \mathcal{L})^2)}
\Big(1+b(p)^2p^{\alpha}
  +b(p)p^{-\frac{1}{2}}(1+p^{\alpha})\Big),
\end{split}
\end{equation}
Choose $\alpha=(\log \mathcal{L})^{-3}$, we see the ratio of the upper bound in \eqref{eqn:Rankinerror} to the right hand side in \eqref{complete sum} is
\begin{equation*}\label{eqn:ratioerror}
\begin{split}
&\ll \exp\Bigg(
  -\alpha\log M+\sum_{\mathcal{L}^2\leq p\leq \exp((\log \mathcal{L})^2)}\Big(
 b(p)^2(p^{\alpha}-1)
 +b(p)p^{-\frac{1}{2}}(p^{\alpha}-1)
  \Big)
  \Bigg)\\
  &\ll \exp\Bigg(
  -\alpha\log M+\sum_{\mathcal{L}^2\leq p\leq \exp((\log \mathcal{L})^2)}\bigg(
  \frac{\mathcal{L}}{p\log p}(\frac{\mathcal{L}}{\log p}+1)(p^\alpha-1)
  \bigg)
  \Bigg)\\
     & \ll \exp(-\alpha\frac{\log M}{\log\log M}).
\end{split}
\end{equation*}
Thus
\begin{equation*}
   |\mathcal{M}_2|\gg
   X\log X \prod_{\mathcal{L}^2\leq p \leq \exp((\log \mathcal{L})^2)}
  \bigg(1+2\frac{b(p)}{\sqrt{p}}+b(p)^2\bigg).
\end{equation*}

Similarly we have
\begin{equation*}
  \mathcal{M}_1\ll X\prod_{\mathcal{L}^2\leq p \leq \exp((\log \mathcal{L})^2)}
\bigg(1+\frac{p b(p)^2}{p+1}\bigg)\ll X\prod_{\mathcal{L}^2\leq p \leq \exp((\log \mathcal{L})^2)}
(1+ b(p)^2).
\end{equation*}
\end{proof}

\end{document}
